\documentclass[12pt,a4paper]{article}

\usepackage[utf8]{inputenc}
\usepackage[T1]{fontenc}
\usepackage{mathpazo}
\usepackage[margin=2.5cm]{geometry}
\usepackage{setspace}
\usepackage{titlesec}
\usepackage[hidelinks,colorlinks=true,linkcolor=black,citecolor=black,urlcolor=blue]{hyperref}
\usepackage{natbib}
\usepackage{url}
\usepackage{footmisc}
\usepackage{csquotes}
\usepackage{abstract}
\usepackage{booktabs}
\usepackage{array}
\usepackage{tabularx}

\onehalfspacing

\titleformat{\section}{\large\bfseries}{\thesection.}{0.5em}{}
\titleformat{\subsection}{\normalsize\bfseries}{\thesubsection.}{0.5em}{}
\titleformat{\subsubsection}{\normalsize\itshape}{\thesubsubsection.}{0.5em}{}

\renewcommand{\abstractnamefont}{\normalfont\bfseries}
\renewcommand{\abstracttextfont}{\normalfont\small}

\title{\textbf{The Industrialization of Plausibility:\\Large Language Models and the Transformation of Semiosis}}

\author{Marco Guerzoni\\[4pt]
\small Department of Economics, Management and Statistics (DEMS)\\
\small Universit\`a degli Studi di Milano-Bicocca\\
\small \url{marco.guerzoni@unimib.it}}
\date{}

\begin{document}

\maketitle

\begin{abstract}
Large Language Models (LLMs) generate discourse that is coherent, contextually
adaptive, and plausible without being grounded in intentional consciousness or
guaranteed reference. This article argues that semiotics---and Umberto Eco's
conceptual vocabulary in particular---provides a uniquely productive framework
for their analysis. It advances three theses: (1)~LLMs instantiate a
computational form of Eco's encyclopedia, organizing meaning through
contextually activated inferential networks rather than hierarchical lemmatic
identity; (2)~they tend structurally toward overcoding and stereotypy,
reproducing culturally dominant frames with an intensity that exceeds even
their training-data distribution; and (3)~their effective use entails a
semiotic cost of context that is unevenly distributed along existing lines of
social stratification. Each thesis is connected to its operational proxy
through an explicit theoretical-empirical bridge that specifies both the
bridging assumption and the condition under which the semiotic interpretation
would be disconfirmed. All three theses are tested through a five-experiment
design moving from mechanism-level disambiguation (E4, E1), through
synchronic and diachronic overcoding measures (E2, E5), to a combined
prompt-factorial and user study (E3). The primary outcome for E2 is a
pre-constructed schema-conformity score, not Shannon entropy; the user study
in E3 measures semiotic cost on the \emph{production} side, not only as a
quality differential. The central claim is that LLMs industrialize the
production of plausible discourse in the specific sense that nineteenth-century
mass production industrialized cultural goods: they separate the production
of form from the labor of meaning, standardize the plausible at scale, and
redistribute onto users and institutions the costs previously borne by the
author.

\smallskip
\noindent\textbf{Keywords:} Large Language Models; semiotics; Umberto Eco;
plausibility; encyclopedia; overcoding; semiotic cost; semiotic inequality;
schema conformity; computational semiotics.
\end{abstract}

\bigskip

%----------------------------------------------------------------------
\section{Introduction}
\label{sec:intro}
%----------------------------------------------------------------------

The emergence of Large Language Models has made visible a condition that
semiotics is uniquely equipped to describe: texts may be coherent, persuasive,
and contextually appropriate without being grounded in either intentional
consciousness or secure reference to the world. LLMs do not ``mean'' in any
human sense, yet they generate utterances that appear meaningful. They do not
``know'' in the robust epistemic sense, yet they produce statements that sound
knowledgeable. What is at stake is not, in the first instance, truth, but
\emph{plausibility}---and plausibility is a semiotic category before it is
an epistemic one.\footnote{Plausibility is here distinguished from
verisimilitude in the Popperian sense \citep{Popper1963}, which concerns
the approximation of a theory to truth. Plausibility designates a property
of discursive form: the capacity to appear acceptable within a given
interpretive framework, independently of verification. It is, in Eco's
terms, a property of the signifying surface rather than a relation to
the world \citep{Eco1976}.}

An utterance appears plausible when it successfully activates culturally
sedimented expectations, recognizable inferential patterns, genre conventions,
and shared background assumptions. It acquires force because it fits an
encyclopedia of possible meanings and acceptable continuations rather than
because it has been tested against the world.

The \emph{industrialization} in the title names a specific transformation of
the productive structure of semiosis, not a simple metaphor for scale. Since
the mid-nineteenth century, industrialization has designated the separation
of production from craft---the extraction of skilled labor's tacit competences
into mechanical processes that can be operated at scale by less-skilled hands.
Applied to discourse, industrialization means the separation of textual
\emph{form} from the communicative \emph{labor} that ordinarily grounds it:
the author's intentional positioning, their lived situatedness, their
accountability to an interpretive community. LLMs achieve exactly this
separation. They produce discourse at scale, at speed, and without the
labor of meaning---that is, without the acts of positioning, commitment,
and accountability that constitute discourse as a social practice. The result
is not meaninglessness but a specifically \emph{industrial} form of plausibility:
standardized, high-volume, and disconnected from the conditions that ordinarily
warrant it.\footnote{This argument intersects with the Frankfurt School's
analysis of the culture industry \citep{AdornoHorkheimer1944}, which identified
in mass cultural production the standardization of form as a mechanism of
ideological reproduction. The present account differs in identifying the
mechanism as semiotic rather than primarily economic, and in focusing on
the redistribution of interpretive labor rather than on the commodification
of aesthetic experience. Eco himself engaged critically with the culture
industry debate in \emph{Apocalittici e integrati} (1964), a text whose
analytical logic is operative here even where it is not cited explicitly.}

A growing body of scholarship has begun to address this phenomenon from semiotic
perspectives. \citet{Leone2023} has outlined the essential tasks of a semiotics
of AI, but does not specify what evidence would distinguish competing semiotic
characterizations. \citet{Vromen2024} has reconceptualized LLMs as semiotic
machines through Saussurean and Derridean theory, but the Derridean framework,
as this article argues, dissolves the distinction between dictionary and
encyclopedic organization rather than enabling its empirical test.
\citet{PiantadosiHill2022} have argued that LLMs capture meaning through
conceptual role, but do not distinguish flat relational structures from
culturally stratified ones. \citet{NunesAntunes2024} have examined the
philosophical challenges of machine semantics without operationalization.
\citet{NotMinds2025} have reframed LLMs through Lotman's semiosphere as
semiotic means rather than cognitive agents---a position this article
endorses and develops empirically. \citet{LSM2025} has argued that LLMs
operate predominantly at the level of the signifier while remaining
disconnected from the conceptual signified---a claim that requires
distinguishing, as the present article does, between dictionary-type and
encyclopedia-type signifier organization. \citet{Tomalin2025} has extended
the analysis to multimodal systems. \citet{HallMirrors2026} has employed
Peircean categories to argue that basic LLMs exist in a hall of mirrors
that newer architectures may be beginning to exit---a Peircean framing
whose relevance to Eco's encyclopedia is discussed below. \citet{Bateman2024}
has offered a discourse-analytic perspective on LLM behavior, providing the
only prior contribution that applies systematic analysis to actual outputs
rather than operating at the level of theoretical characterization.
\citet{LackovaBennett2025} has proposed a biosemiotic framework connecting
folding dynamics to meaning-production.\footnote{See also \citet{Signata2025}
for a recent thematic issue on AI and digital semiotic spaces.}

What is common to all of these contributions, with the partial exception of
\citet{Bateman2024}, is that their central claims remain at the level of
theoretical characterization: they propose frameworks, draw analogies, and
identify challenges, but they do not specify the conditions under which
their characterizations would be \emph{wrong}. An analogy that cannot fail
is not analytically productive, and an account of LLM semiosis that fits
any possible LLM behavior is not explanatory. This article addresses the
gap by formulating each theoretical claim as a falsifiable prediction,
specifying the operational proxy through which it is measured, and testing
it against actual LLM outputs with pre-registered analysis plans.

The argument is structured as follows. Section~\ref{sec:framework} introduces
the three Eco concepts that ground the analysis. Section~\ref{sec:theses}
formulates the three theses with explicit forward pointers to their experimental
tests. Section~\ref{sec:bridge} formalizes the theoretical-empirical bridge.
Section~\ref{sec:experiment} presents the five-experiment design.
Section~\ref{sec:implications} derives the conditional implications.
Section~\ref{sec:discussion} addresses the theory-measurement gap, scope
limits, and adjacent literatures. Section~\ref{sec:conclusion} concludes.


%----------------------------------------------------------------------
\section{Eco's Semiotics as Analytical Framework}
\label{sec:framework}
%----------------------------------------------------------------------

LLMs do not operate on things but on signifying forms. Their domain is
textuality, discursivity, codification, and patterned relations among linguistic
units. Even when they produce statements about the world, they do so through
structures internal to discourse. This is why semiotics is not an ornamental
perspective on LLMs but one of the most suitable disciplines for their analysis:
it studies the conditions under which something functions as a sign, and LLMs
are entirely dependent on those conditions.\footnote{This situates the present
approach within what \citet{Harnad1990} calls the symbol grounding problem:
the question of how symbols acquire meaning that is not merely parasitic on
other symbols. LLMs, as their critics have noted (e.g., \citealt{Searle1980}),
operate on ungrounded symbols---their ``meaning'' is derived, not original, in
Harnad's sense. The present article does not resolve this debate but operates
at a different level: not whether LLM signs are genuinely meaningful, but how
the semiotic properties of their outputs function socially and culturally,
regardless of their grounding status. This is a legitimate and important
question even if grounding is denied.}

Among semiotic traditions, Eco's provides three concepts of particular
analytical power for the study of LLMs: the \emph{encyclopedia},
\emph{overcoding}, and \emph{interpretive cooperation}. These are not merely
suggestive metaphors; they are theoretically precise tools that, as the
experimental section shows, can be operationalized and tested.

\subsection{The Encyclopedia}

For Eco, meaning is not exhausted by dictionary definition. It depends on
an open, distributed network of cultural knowledge, interpretive habits,
practical scripts, intertextual memories, and socially shared assumptions
\citep{Eco1984}. The dictionary model treats meaning as a fixed set of
necessary and sufficient features organized in a taxonomic tree. The
encyclopedia model treats meaning as a potentially unlimited network of
culturally organized interpretants, navigable in different directions
depending on contextual cues \citep{Eco1976,Eco1984}.\footnote{Eco's
encyclopedia model transforms Peirce's concept of unlimited semiosis, in
which every interpretant becomes a sign requiring further interpretation
\citep{Peirce1931}. Unlike Peirce, Eco insists on the pragmatic and
cultural constraints that delimit this process in practice \citep{Eco1990}.
The Peircean framing matters for LLMs: \citet{HallMirrors2026} has argued
that basic LLMs are trapped in a chain of signifier-to-signifier relations
with no thirdness---no interpretant that connects the sign to action or
habit. The encyclopedia thesis adopted here makes a weaker claim: not that
LLMs have genuine interpretants, but that their outputs are organized by the
same culturally weighted pathways that structure human encyclopedic access.
This is a claim about the functional organization of outputs, not about
grounding.}

The distinction is analytically precise: in a dictionary, ``crane'' has a
primary meaning (a bird) and secondary meanings (a machine, a verb) organized
hierarchically. In an encyclopedia, ``crane'' is a node connected to
ornithology, construction, Japanese art, port logistics, origami, and Ichabod
Crane through culturally weighted but non-hierarchical pathways. Which pathway
is activated depends on context.

This distinction matters for LLMs because it generates two testable predictions.
The first is \emph{mechanistic}: a system with encyclopedic organization should
disambiguate polysemous items in ways that require multi-hop cultural inference
unavailable to a frequency-ranked sense dictionary---and crucially, it should
do so on items where encyclopedic inference and statistical co-occurrence
predict \emph{divergent} responses. The second is \emph{distributional}: at
scale, context-based grouping of outputs should produce more coherent clusters
than lemma-based grouping, and this effect should persist in conditions where
domain vocabulary is absent from the prompt.

\subsection{Overcoding}

Eco uses the term \emph{overcoding} to describe the superimposition of
additional cultural codes upon a sign or text, constraining its interpretation
toward culturally stabilized pathways \citep{Eco1976}.\footnote{Overcoding
refers to the rule-governed process by which a sign already codified at one
level receives further cultural specification at a higher level, directing
interpretation toward the most conventional reading. It is not a defect but
a structural property of any signifying practice that has accumulated social
convention. See \citet{Eco1976}, Chapter~2.} When overcoding is pervasive,
meaning follows paths of least resistance: the most conventional, the most
familiar, the most culturally entrenched.

LLMs are, by construction, overcoding machines. Their generative mechanism
privileges high-probability continuations---tokens most frequently associated
with a given context in the training corpus. The most probable continuation
is, by definition, the most culturally sedimented one. This produces outputs
that conform to dominant narrative schemas, rhetorical conventions, and
ideological defaults---not because the machine ``believes'' them, but because
they constitute the path of least statistical resistance.

Overcoding, in Eco's account, is not merely a synchronic property of a code
but a social-historical process: conventions solidify as cultural usage
accumulates and pathways become more deeply grooved. This diachronic dimension
connects directly to Lotman's theory of the semiosphere \citep{Lotman1990}.
For Lotman, semiotic space is organized around a center where codes are most
stable, most explicit, and most self-describing, and a periphery where codes
are less stabilized and more open to transformation. The semiosphere does not
evolve uniformly: it undergoes periods of gradual sedimentation punctuated by
explosive reorganization when peripheral formations break through into the
center. LLMs, trained predominantly on majority-language, majority-culture
corpora, implicitly model the center of specific semiospheres at specific
historical moments. Their outputs therefore reinforce center-periphery
asymmetries not merely by reproducing them, but by removing the creative
friction---misunderstanding, idiolectal deviation, translational loss---that,
in Lotman's account, is the primary mechanism of semiospheric renewal. An
encyclopedia that learns from a corpus, then generates text that is fed back
into subsequent training corpora, is not merely a passive recorder of the
semiosphere: it is a mechanism for accelerating sedimentation at the center
while starving the periphery of the generative noise that drives semiotic
evolution.

The overcoding thesis thus generates two empirical predictions. Synchronically,
outputs for culturally charged concepts should \emph{conform to pre-identified
cultural schema templates} more closely than outputs for matched culturally
neutral concepts---and this schema conformity should exceed the conformity
already present in the training corpus. Diachronically, schema conformity
should increase across successive model generations for concepts of rising
cultural salience, while remaining stable for concepts whose cultural charge
has not changed.

\subsection{Interpretive Cooperation and the Pragmatics of Context}

Eco argues that texts are designed to activate a Model Reader capable of
actualizing their semantic potential \citep{Eco1979}. A text does not fully
determine its own meaning; it requires a competent interpreter who activates
the relevant codes, supplies background knowledge, and restrains arbitrary
readings \citep{Eco1979,Eco1990}. Meaning is therefore a cooperative
achievement, not a unilateral emission.

LLM-generated discourse depends even more radically on such cooperation.
The user must specify the frame, delimit the task, supplement background
assumptions, and judge the output's adequacy. But this cooperation is
asymmetric in a distinctive way: what remains implicit in human communication---
shared institutional context, mutual knowledge, situational awareness---must
often be made \emph{explicit} in the prompt. Context, which is ordinarily a
background condition of interaction, becomes a designed input that requires
labor.

This labor is what Eco's concept of the Model Reader illuminates and what
Bourdieu's analysis of linguistic capital explains sociologically
\citep{Bourdieu1991}. Eco's Model Reader is a \emph{textual} construct: it
names the set of competences that a text presupposes in its ideal interpreter.
But a text's Model Reader is not distributed randomly across the social
space. Bourdieu's linguistic capital---the ensemble of linguistic resources
whose possession confers communicative authority and access to legitimate
discourse---is precisely the social mechanism by which the competences
presupposed by authoritative texts are distributed unequally along lines
of class, education, and cultural origin. The connection between the two
concepts is therefore not a juxtaposition but a theoretical articulation:
Eco's Model Reader specifies \emph{what} competences are required; Bourdieu's
linguistic capital explains \emph{why} those competences are not equally
available. For LLMs, this articulation generates the concept of
\emph{semiotic cost of context}: the labor required to construct the
interpretive frame that activates the model's relevant encyclopedic pathways
is a new form of the old asymmetry.

This cost is measurable on both sides of the exchange. On the output side,
the quality of LLM-generated discourse varies with the contextual richness
of the prompt---the model activates better encyclopedic pathways when given
more precise framing. On the production side, constructing a precise frame
costs labor---time, cognitive effort, iterative refinement, failed attempts,
and the frustration of receiving outputs that are not what was needed. The
inequality is not only that some users receive better outputs; it is that
some users expend more labor to receive worse ones.


%----------------------------------------------------------------------
\section{Three Theses}
\label{sec:theses}
%----------------------------------------------------------------------

\subsection{Thesis~1: LLMs Instantiate a Computational Form of the Encyclopedia}
\label{sec:thesis1}

LLMs do not operate as static lexicons. They approximate networks of
associations, semantic neighborhoods, discursive regularities, and latent
inferential paths. Their outputs emerge from encyclopedia-like organization
rather than from dictionary-like storage \citep{Eco1984}.

This claim requires precise formulation. The analogy between Eco's encyclopedia
and LLM latent spaces can be drawn at different levels of rigor. At the weakest
level, it is merely suggestive: LLMs ``associate'' words in ways that
``resemble'' cultural knowledge. This level covers most of the existing
literature. \citet{Vromen2024} treats the LLM's latent space as a model of
Derridean \emph{\'ecriture}, where meaning arises from differential relations
among signifiers. \citet{LSM2025} has corroborated the claim from a Saussurean
perspective, noting that LLMs operate predominantly at the level of the
signifier. Both accounts are illuminating, but neither specifies what it would
mean for the analogy to \emph{fail}. An analogy that cannot fail is not
analytically productive.

At a stronger level, the encyclopedic thesis generates falsifiable predictions.
But there is a specification problem: superior disambiguation in context-rich
conditions is predicted by \emph{any} sufficiently large language model with
adequate context window, not only by one with encyclopedic organization in Eco's
sense. A model that performs well-calibrated co-occurrence statistics at scale
would produce identical results to a model with genuine cultural-encyclopedic
organization, in standard disambiguation tasks. The critical empirical test is
therefore not whether context helps---it does, trivially---but whether LLM
disambiguation is \emph{governed by encyclopedic inference in cases where
encyclopedic inference and high-frequency co-occurrence predict divergent
outcomes}.

This requires stimuli with two properties: (a) a polysemous item where one
reading is statistically dominant in the training corpus, and (b) a contextual
frame where the culturally appropriate reading is the \emph{non-dominant} one,
requiring cultural schema activation rather than frequency-weighted sense
assignment. If LLMs systematically choose the encyclopedic reading on these
items---outperforming both the dictionary model and a co-occurrence-based
baseline---the encyclopedia thesis is empirically supported in a way that
survives the frequency-confound objection. This design principle governs
Experiment~E4 (Section~\ref{sec:exp4}).

The distributional extension of the thesis (Experiment~E1) addresses a
complementary question: across a broad lexical sample, does context-based
clustering of outputs exceed lemma-based clustering, and does this hold even
when domain vocabulary is absent from the prompt? Confirmation of E4 and E1
provides \emph{independent, converging evidence} for the encyclopedic
characterization. Each is informative alone; together they constitute a
two-level confirmation. Disconfirmation of E4 alone would falsify the
mechanistic prediction; disconfirmation of E1 alone would leave the
mechanistic prediction untouched while falsifying the distributional one.
Full disconfirmation of Thesis~1 requires failure of \emph{both} predictions.

\subsection{Thesis~2: LLMs Tend Structurally Toward Overcoding and Stereotypy}
\label{sec:thesis2}

Because LLMs privilege high-probability continuations, they tend to reproduce
already stabilized codifications. They are structurally inclined toward what
is discursively dominant.

The existing literature has addressed this tendency primarily under the rubric
of ``bias'' \citep{Bender2021}. But overcoding identifies something broader and
more fundamental than bias in the computational fairness sense. Bias implies a
correctable deviation from a norm of neutrality. Overcoding identifies a
structural property: the most statistically salient path through a culturally
sedimented corpus \emph{is} the stereotype. There is no neutral distribution to
return to; there is only the culturally weighted one. Corrective interventions
must therefore either restructure the encyclopedia (modifying whose conventions
are dominant) or impose post-hoc constraints that substitute one overcoding
for another.

This structural claim requires an operational proxy that measures not just
distributional variance across outputs (which entropy measures) but conformity
to \emph{pre-identified cultural schema templates}. Shannon entropy is an
insufficient proxy for overcoding because it is agnostic about the content
of the frames it counts: a low-entropy distribution of \emph{diverse} frames
would look identical to a low-entropy distribution of \emph{stereotypic} ones.
What overcoding predicts is not simply convergence but convergence onto a
specific culturally dominant form. The primary outcome for Experiment~E2 is
therefore a schema-conformity score: the degree to which each output matches
a pre-constructed cultural prototype, independently specified by domain
experts before data collection. Entropy is retained as a secondary measure
of distributional concentration.

The diachronic prediction concerns not the mechanism driving overcoding across
model generations (which cannot be isolated from architectural and RLHF
differences across generations) but the observable trajectory: schema
conformity should increase across successive model generations for concepts
of rising cultural salience, and remain stable for stable-salience controls.
This is a weaker but empirically tractable claim, tested by Experiment~E5.

\subsection{Thesis~3: The Construction of Context Entails a Semiotic Cost
That Is Unevenly Distributed}
\label{sec:thesis3}

What remains implicit in human interaction often must be made explicit for the
model. To obtain relevant output, users must specify aims, roles, assumptions,
genres, constraints, and background knowledge. Context therefore requires labor,
and this labor is unevenly distributed along existing lines of social
stratification.

This thesis requires evidence of asymmetry on \emph{both sides} of the
exchange. The output side is the easier half: it can be shown that prompt
sophistication predicts output quality. The harder and theoretically more
important half is the production side: that users with less semiotic capital
invest more labor to receive worse outputs. A study that only measures output
quality differences across prompt levels demonstrates that sophisticated
prompts yield better results---which is unsurprising and says nothing about
inequality. The semiotic inequality claim requires demonstrating that the
labor cost of producing a sophisticated prompt is itself unevenly distributed:
that users with lower semiotic competence spend more effort, iterate more
often, experience more failure, and arrive at prompts of lower sophistication
despite the effort investment.

Experiment~E3 (Section~\ref{sec:exp3}) is designed to address both sides.
The prompt-factorial component tests the output side. The user study
($N = 120$, stratified by education level and domain familiarity) measures
the production side by recording time-on-task, iteration count, self-reported
difficulty, and component sophistication of spontaneously produced prompts.
The pre-registered inequality claim requires that prompt sophistication is
positively correlated with output quality \emph{and} negatively correlated
with production effort---so that users with higher semiotic competence spend
less effort to produce better prompts that yield better outputs. Users with
lower competence spend more effort, produce less sophisticated prompts, and
receive worse outputs: this is the double bind that constitutes semiotic
inequality in its specifically computational form.


%----------------------------------------------------------------------
\section{Theoretical-Empirical Bridge}
\label{sec:bridge}
%----------------------------------------------------------------------

Before presenting the experimental design, it is necessary to make explicit
the bridging assumptions that connect Eco's theoretical constructs to the
computational proxies employed in the experiments. Semiotic theory operates
at the level of social meaning-making processes; the experiments operate at
the level of measurable properties of LLM outputs. These levels are not
identical. The validity of the empirical claims depends on whether the
proxies track the theoretical constructs with sufficient fidelity, and the
theoretical claims are evaluable only if the bridging assumptions are explicit
enough to be challenged.

Table~\ref{tab:bridge} formalizes these connections for all five experiments.
Three features of the table deserve comment.

\begin{table}[ht]
\centering
\small
\renewcommand{\arraystretch}{1.55}
\begin{tabularx}{\textwidth}{p{2.2cm}p{2.8cm}p{3.2cm}p{3.5cm}}
\toprule
\textbf{Eco Construct} & \textbf{Theoretical Definition} &
\textbf{Operational Proxy (Experiment)} &
\textbf{Falsifiability Condition} \\
\midrule
Encyclopedia &
Network of culturally organized interpretants navigated by contextual
activation, not lemmatic hierarchy &
(E4) LLM accuracy on divergent-prediction critical items exceeds co-occurrence
baseline; (E1) Context-based silhouette $>$ lemma-based, replicating in
paraphrase-context condition &
E4: LLMs fail to exceed co-occurrence baseline on critical items with frequency-conflicting
encyclopedic readings. E1: Context silhouette $\leq$ lemma silhouette, or
paraphrase condition fails to replicate. Either alone partially disconfirms;
both together fully disconfirm Thesis~1 \\
\addlinespace
Overcoding &
Superimposition of culturally stabilized codes constraining interpretation
toward dominant convention; intensifies diachronically &
(E2) Schema-conformity scores for culturally charged concepts exceed matched
controls and training-corpus baseline; (E5) Schema conformity increases
across model generations for rising-salience concepts, stable for controls &
E2: Charged concepts do not exceed matched controls in schema conformity, or
LLM scores do not exceed corpus baseline. E5: Generation $\times$
concept-type interaction non-significant. Either alone partially disconfirms;
both together fully disconfirm Thesis~2 \\
\addlinespace
Semiotic cost of context &
Labor of constructing an interpretive frame for machine-generated discourse;
unevenly distributed along lines of semiotic competence &
(E3) Expert-verified output quality increases with contextual richness (output
side); production effort (time, iterations, component sophistication) is
negatively correlated with semiotic competence after controlling for domain
knowledge (production side) &
If output quality does not increase monotonically, or if the domain
$\times$ level interaction is non-significant, the output-side prediction
fails. If production effort is uncorrelated with semiotic competence
after controls, the inequality claim fails. Both must hold for full
confirmation of Thesis~3 \\
\bottomrule
\end{tabularx}
\caption{Theoretical-empirical bridge: constructs, operational proxies,
and falsifiability conditions. E1--E5 refer to the five experiments in
Section~\ref{sec:experiment}. Confirmation and disconfirmation logic
for each thesis is specified in the rightmost column.}
\label{tab:bridge}
\end{table}

\paragraph{On the E4 baseline design.} The critical innovation in E4 is the
inclusion of \emph{divergent-prediction critical items}: stimuli for which the
statistically most frequent sense and the encyclopedically appropriate sense
differ. On these items, a pure co-occurrence model will systematically choose
the wrong reading; a system with encyclopedic cultural organization will choose
correctly. Superior LLM performance on exactly these items is the specific
prediction that the co-occurrence-confound objection cannot explain away.

\paragraph{On schema conformity vs.\ entropy (E2).} Shannon entropy measures
how concentrated a distribution is, not what it is concentrated \emph{on}.
Two distributions could have identical entropy while one clusters on a
culturally stereotypic frame and the other on a non-stereotypic one. Schema
conformity, operationalized as the mean cosine similarity between output
embeddings and a pre-constructed prototype embedding for each concept
(specified before data collection by domain experts), measures what overcoding
predicts: not mere convergence but convergence onto a culturally dominant form.
Entropy is retained as a secondary measure; the two measures are expected to
correlate but are conceptually distinct.

\paragraph{On multi-model replication.} All five experiments are run on three
architecturally distinct models: GPT-4o, Llama-3-70B-Instruct, and
Mistral-Large-2. These models differ in training corpus, architecture scale,
instruction-tuning procedure, and RLHF implementation. Claims are restricted
to patterns that replicate across all three. Where results diverge, the
divergence is analyzed with reference to known architectural differences;
model-specific results are labeled as such and not presented as properties of
LLMs as a class. The three-model replication is the minimum condition for
treating experimental results as evidence about structural tendencies of
instruction-tuned LLMs rather than about idiosyncrasies of a particular
system.


%----------------------------------------------------------------------
\section{Experimental Design}
\label{sec:experiment}
%----------------------------------------------------------------------

The five experiments are ordered by logical dependency. E4 and E1 establish
the encyclopedic mechanism. E2 and E5 establish the overcoding tendency
synchronically and diachronically. E3 tests the social-access thesis.
Table~\ref{tab:exp_summary} summarizes the design. All study materials,
pre-registration documents, and analysis code will be made available on
OSF prior to data collection (pre-registration DOI to be inserted upon
filing).

\subsection{Shared Methods}
\label{sec:shared}

\paragraph{Models.} GPT-4o (\texttt{gpt-4o-2024-08-06}), Llama-3-70B-Instruct
(Together~AI), Mistral-Large-2 (Mistral~AI API). Model versions are fixed and
logged.

\paragraph{Prompting.} System prompt constant across all conditions:
``You are a helpful assistant. Respond in English only.'' Experimental
manipulations are introduced exclusively through the user turn.

\paragraph{Reproducibility.} All API calls are logged (timestamp, model,
prompt ID, temperature, seed, response, token count). UMAP hyperparameters
are fixed by grid search on a held-out development set before main analysis.
Random seeds are fixed and reported.

\paragraph{Embeddings.} Two embedding models from different providers to
avoid co-training circularity: \texttt{text-embedding-3-large} (OpenAI)
and \texttt{e5-large-v2} \citep{Wang2022} (Microsoft/Hugging Face).
Claims restricted to patterns replicating under both.

\paragraph{Pre-registration.} All hypotheses, analysis plans, and stopping
rules pre-registered on OSF before data collection. The pre-registration
document specifies: the directional prediction for each experiment, the
analysis model (including random effects structure), criteria for accepting
the LLM-as-judge in E3, and the adjudication rule for conflicting results
within a thesis (see below).

\paragraph{Adjudication rules.} Where experiments within the same thesis
yield contradictory results (one confirms, one disconfirms), the adjudication
rule pre-registered for each thesis governs the interpretation:
\begin{itemize}
\item \emph{Thesis~1}: E4 (mechanistic) takes precedence over E1
      (distributional); E4 confirmation with E1 disconfirmation supports a
      restricted version of the thesis (mechanistic but not distributional).
\item \emph{Thesis~2}: E2 (synchronic) and E5 (diachronic) are independent;
      either may be confirmed or disconfirmed independently; both must be
      confirmed for the full diachronic claim to stand.
\item \emph{Thesis~3}: The output-side (prompt-factorial) and production-side
      (user study) predictions are necessary for the \emph{inequality} claim
      but the output-side is sufficient for a weaker quality-differential claim.
\end{itemize}

\begin{table}[ht]
\centering
\small
\renewcommand{\arraystretch}{1.4}
\begin{tabularx}{\textwidth}{p{0.7cm}p{1.5cm}p{1.7cm}p{2.8cm}p{2.8cm}p{2.5cm}}
\toprule
\textbf{Exp.} & \textbf{Thesis} & \textbf{Outputs} & \textbf{Key analysis} &
\textbf{Confound addressed} & \textbf{Primary outcome} \\
\midrule
E4 & T1 & 6,000 &
Disambiguation accuracy by system $\times$ depth; critical items analysis &
Frequency confound; prompt injection &
Accuracy on divergent-prediction items \\
\addlinespace
E1 & T1 & 9,000 &
Silhouette (context vs.\ lemma); paraphrase condition; baselines &
Embedder circularity; lexical priming &
Silhouette differential, paraphrase condition \\
\addlinespace
E2 & T2 & 48,000 &
Schema conformity vs.\ corpus baseline; entropy secondary &
Entropy agnostic to frame content; corpus confound &
Schema-conformity score \\
\addlinespace
E5 & T2 & 4,000 &
Conformity trajectories across 4 generations &
Synchronic artifact &
Generation $\times$ concept-type interaction \\
\addlinespace
E3 & T3 & 14,400+ &
Mixed-effects quality model; user study ($N=120$) double bind &
LLM-judge circularity; output-only operationalization &
Quality (output) and effort-to-quality ratio (production) \\
\bottomrule
\end{tabularx}
\caption{Experimental overview: thesis, outputs, key analysis, confound
addressed, and primary outcome for each experiment.}
\label{tab:exp_summary}
\end{table}

%----------------------------------------------------------------------
\subsection{Experiment~E4: Dictionary--Encyclopedia Diagnostic}
\label{sec:exp4}
%----------------------------------------------------------------------

\emph{Tests Thesis~1 (mechanistic prediction).}

\subsubsection{Rationale}

The encyclopedic thesis claims that LLMs use cultural-inferential context,
not frequency-ranked sense assignment. This claim is testable only on stimuli
where the two mechanisms predict \emph{different} outcomes. Standard
disambiguation tasks do not satisfy this requirement because a sufficiently
capable co-occurrence model will outperform a dictionary model and coincidentally
agree with the encyclopedic prediction on most items. The design must include a
critical-items stratum---items where the most frequent sense in the training
corpus is \emph{not} the contextually appropriate one---so that statistical
advantage and encyclopedic accuracy can be dissociated.

\subsubsection{Materials}

Fifty minimal-pair sentences are constructed. Each pair consists of:
(a) a sentence activating the statistically dominant sense of the target item,
and (b) a sentence activating a culturally appropriate but frequency-\emph{non-}dominant
sense through inferential context. The pairs are stratified across three
inferential-depth categories:

\begin{itemize}
\item \textbf{Shallow} ($n = 15$): single-hop cultural inference
      (e.g., ``the crane at the construction site'' vs.\ ``the crane by the
      reed beds'').
\item \textbf{Medium} ($n = 15$): two-hop script knowledge
      (e.g., a sentence where the appropriate reading requires activating a
      professional script that is culturally salient but corpus-rare).
\item \textbf{Deep} ($n = 20$): requires culturally specific schema activation
      whose co-occurrence trace in general corpora is low.
\end{itemize}

\noindent\textbf{Critical items.} Within each stratum, a subset of items
(target: $n = 8$ per stratum, $n = 24$ total) is designated as
\emph{critical items}: sentences of type~(b) where corpus-frequency analysis
confirms that the encyclopedic reading is the non-dominant one in the training
distribution. This is verified using WordNet sense frequency data
\citep{Miller1995} and sense-tagged corpus statistics. On these items, a
co-occurrence model predicts the wrong reading; an encyclopedic model predicts
the correct one. The critical-items analysis is the primary test of Thesis~1.

Ground truth labels are assigned by five independent annotators (professional
linguists, lexical semantics specialization); Fleiss' $\kappa \geq 0.80$
required; items below threshold are replaced.

\subsubsection{Procedure}

Each sentence is presented with the probe: ``What does the word `[TARGET]' refer
to in the following sentence? Provide a one-sentence explanation.'' $k = 20$
responses per sentence per model at $T = 0.0$ (deterministic; primary analysis)
and $T = 0.7$ (distributional analysis), yielding $50 \times 20 \times 3 = 3{,}000$
responses per temperature condition. Responses are classified (encyclopedic /
dictionary / neither) by a rubric-based classifier validated on a held-out
set (Cohen's $\kappa$ reported).

\subsubsection{Baselines}

(1)~\emph{Dictionary model}: most frequent WordNet sense. (2)~\emph{WSD
system}: \citet{Bevilacqua2020} \texttt{EWISER}, trained on SemEval data.
(3)~\emph{Co-occurrence model}: a large n-gram language model with no
instruction tuning (GPT-2 XL) used purely as a frequency-weighted sense
predictor. (4)~\emph{Human ceiling}: na\"ive annotators, $n = 30$.

\subsubsection{Analysis}

Disambiguation accuracy is computed per system per inferential-depth stratum,
with a separate analysis restricted to critical items. A mixed-effects logistic
regression is fitted: correct $\sim$ system $\times$ depth $\times$ item\_type
(critical vs.\ non-critical) + (1$|$sentence) + (1$|$lemma). The pre-registered
primary test is the three-way interaction: LLMs should outperform the
co-occurrence baseline specifically on critical items at medium and deep
inferential depth.

%----------------------------------------------------------------------
\subsection{Experiment~E1: Encyclopedic Organization at Distributional Scale}
\label{sec:exp1}
%----------------------------------------------------------------------

\emph{Tests Thesis~1 (distributional prediction).}

\subsubsection{Materials}

Thirty polysemous terms are drawn from a stratified OED sample:
\emph{high-ambiguity} ($n = 12$; e.g., \emph{bank}, \emph{crane},
\emph{pitch}, \emph{cell}), \emph{medium-ambiguity} ($n = 12$), and
\emph{near-monosemous controls} ($n = 6$). For each term, five prompt
variants are constructed:

\begin{enumerate}
\item \textbf{Neutral}: ``Tell me about [term].''
\item \textbf{Explicit-A}: Domain named explicitly (dominant sense).
\item \textbf{Explicit-B}: Domain named explicitly (secondary sense).
\item \textbf{Paraphrase-A}: Domain implied through narrative framing
      without domain vocabulary (dominant sense).
\item \textbf{Paraphrase-B}: Domain implied through narrative framing
      without domain vocabulary (secondary sense).
\end{enumerate}

Paraphrase conditions (4--5) are critical: context-sensitivity that persists
without domain vocabulary in the prompt cannot be attributed to lexical priming.

\subsubsection{Procedure}

$30 \times 5 \times 20 \times 3 = 9{,}000$ outputs at $T = 0.7$, embedded under
both embedding models (Section~\ref{sec:shared}).

\subsubsection{Analysis}

For each of the six embedding corpora (3 models $\times$ 2 embedders):

\begin{enumerate}
\item UMAP dimensionality reduction with hyperparameters fixed by grid
      search on a held-out development set. \emph{Stability protocol}:
      five independent UMAP runs per corpus; silhouette scores are averaged
      across runs; items with cross-run variance exceeding a pre-specified
      threshold ($\text{SD} > 0.05$) are flagged.
\item Two clustering analyses: by lemma and by context. Silhouette scores
      per grouping, stratified by ambiguity stratum.
\item Pre-registered: silhouette(context) $>$ silhouette(lemma) in high- and
      medium-ambiguity strata; silhouette(lemma) $\geq$ silhouette(context) in
      monosemous controls.
\item Paraphrase-specific analysis: the silhouette difference
      silhouette(context) $-$ silhouette(lemma) is computed separately for
      explicit conditions (2--3) and paraphrase conditions (4--5). A significant
      positive difference in paraphrase conditions rules out lexical priming.
\end{enumerate}

Three baselines: TF-IDF, GloVe \citep{Pennington2014}, outputs of
\texttt{EWISER} \citep{Bevilacqua2020}. LLM context-based silhouette scores
must exceed all three baselines for the encyclopedic claim.

%----------------------------------------------------------------------
\subsection{Experiment~E2: Schema Conformity and Overcoding}
\label{sec:exp2}
%----------------------------------------------------------------------

\emph{Tests Thesis~2 (synchronic prediction).}

\subsubsection{Rationale}

Overcoding predicts conformity to a dominant cultural frame, not merely
low output variance. The primary outcome is therefore schema conformity:
the degree to which each output matches a pre-constructed cultural prototype
for each concept. Shannon entropy over a topic-model distribution is a
secondary measure of distributional concentration. The two measures address
distinct questions: schema conformity asks ``does the output conform to the
dominant cultural frame?''; entropy asks ``how diverse are the frames
generated?'' Both are relevant; neither alone is sufficient.

\subsubsection{Schema prototype construction}

For each of the 20 culturally charged concepts, a prototype embedding is
constructed before any LLM data collection. Prototypes are derived from:
(a) a structured elicitation task administered to $n = 30$ culturally
representative human raters (who list the five features most associated
with the concept), (b) the centroid of the top-10 Wikipedia passages most
frequently associated with the concept, and (c) existing cultural stereotype
research for person-concepts (e.g., \citealt{Fiske2002} for
scientist/leader/criminal). The prototype embedding is the sentence embedding
of the synthesized prototype description; all embeddings use
\texttt{e5-large-v2} to avoid circularity with the generative models.

\subsubsection{Materials}

Twenty culturally charged concepts (selected via Schwartz Basic Human Values
taxonomy; cultural stereotypy rated by 5 annotators, ICC $\geq 0.70$) and 20
matched controls (COCA frequency $\pm 15\%$, Wikipedia length $\pm 20\%$,
stereotypy $< 3.0/7.0$). Prompt: ``Write a short paragraph describing
[concept].'' Corpus baselines: 50 Wikipedia paragraphs + 100 Common Crawl
documents per concept, embedded identically.

\subsubsection{Procedure}

For each of 40 concepts $\times$ 3 models $\times$ 4 temperatures
$T \in \{0.2, 0.5, 0.8, 1.2\}$: $k = 100$ responses $= 48{,}000$ outputs.

\subsubsection{Analysis}

\paragraph{Schema conformity (primary).} Mean cosine similarity between each
output's embedding and its concept's prototype embedding, per concept per
temperature per model. A mixed-effects model is fitted: schema\_conformity
$\sim$ concept\_type (charged vs.\ control) $\times$ temperature +
(1$|$concept). Pre-registered primary test: charged concepts exhibit
significantly higher schema conformity than matched controls, and LLM
schema conformity for charged concepts exceeds corpus-baseline conformity.
The latter comparison is the model-amplification test.

\paragraph{Entropy (secondary).} Shannon entropy $H = -\sum_i p_i \log_2 p_i$
is computed from BERTopic distributions (three hyperparameter configurations;
Spearman $\rho \geq 0.85$ cross-configuration stability required).
The adjudication rule for conflicting primary and secondary measures: if
schema conformity is high but entropy is also high (multiple stereotypic
frames), the overcoding thesis is confirmed; if schema conformity is low
but entropy is also low (convergence on a non-stereotypic frame), the result
is considered anomalous and requires case-by-case analysis. Schema conformity
takes precedence as the theoretically grounded primary outcome.

\paragraph{Temperature ablation.} Piecewise linear regression of schema
conformity on temperature tests for a conformity ceiling: the prediction is
that schema conformity does not decrease monotonically with temperature for
charged concepts (structural floor), while it increases monotonically for
neutral controls (no structural constraint).

\paragraph{Demographic markers (person-concepts).} For six person-concepts
(\emph{scientist}, \emph{leader}, \emph{criminal}, \emph{nurse},
\emph{professor}, \emph{entrepreneur}): gender, age, and ethnicity markers
are extracted via a calibrated classifier and compared against corpus
baselines via McNemar's test. This sub-analysis is reported descriptively
given the small number of concepts ($n = 6$); no omnibus inference is drawn.

%----------------------------------------------------------------------
\subsection{Experiment~E5: Diachronic Trajectory of Schema Conformity}
\label{sec:exp5}
%----------------------------------------------------------------------

\emph{Tests Thesis~2 (diachronic prediction).}

\subsubsection{Rationale}

The diachronic prediction concerns the observable trajectory of schema
conformity across model generations, not the mechanism that drives it.
Model generations differ in training corpus, architecture, and RLHF
procedure simultaneously; the design does not attempt to isolate any single
cause. The pre-registered claim is therefore descriptive-diachronic:
schema conformity increases across generations for rising-salience concepts
and remains stable for stable-salience controls.

\subsubsection{Materials}

Ten \emph{changing-salience} concepts (Google Trends volume and Wikipedia
edit frequency increased significantly between 2019 and 2024; e.g.,
\emph{pandemic}, \emph{remote work}, \emph{artificial intelligence},
\emph{vaccine}, \emph{disinformation}) and ten \emph{stable-salience}
controls. Prototype embeddings from E2 are reused where applicable; for
concepts not in the E2 set, new prototypes are constructed by the same
procedure.

\subsubsection{Procedure}

E2 protocol ($k = 50$, $T = 0.8$, BERTopic + schema conformity) on four
model generations: GPT-3 (\texttt{text-davinci-001}), GPT-3.5
(\texttt{gpt-3.5-turbo-0125}), GPT-4 (\texttt{gpt-4-0613}), GPT-4o. Where
legacy endpoints are unavailable, archival documentation is provided.
Total outputs: $20 \times 50 \times 4 = 4{,}000$.

\subsubsection{Analysis}

Schema conformity trajectories are plotted per concept type. Mixed-effects
model: schema\_conformity $\sim$ generation (ordinal) $\times$
concept\_type + (1$|$concept). Pre-registered critical test: the interaction
term. A significant positive interaction (conformity increases across
generations for changing-salience concepts, not for controls) confirms the
diachronic prediction.

%----------------------------------------------------------------------
\subsection{Experiment~E3: The Semiotic Cost of Context}
\label{sec:exp3}
%----------------------------------------------------------------------

\emph{Tests Thesis~3.}

\subsubsection{Rationale}

The semiotic inequality thesis requires evidence on \emph{both sides} of the
exchange: output quality should increase with prompt sophistication (output
side), and users with lower semiotic capital should invest more effort to
produce less sophisticated prompts and receive worse outputs (production side).
The user study is designed to operationalize the double-bind structure of
semiotic inequality, not merely to demonstrate that prompt quality varies.

\subsubsection{Materials}

\paragraph{Domains.} Ten task domains pre-classified by independent expert
ratings into high-complexity ($n = 5$: landlord-tenant legal dispute, medical
test interpretation, pharmaceutical adverse event report, financial instrument
evaluation, regulatory compliance) and standard-complexity ($n = 5$: rental
agreement review, dietary recommendation, job application cover letter,
insurance claim, academic appeal letter). ICC $> 0.75$ required.
Expert-verified ground-truth benchmark responses constructed per domain.

\paragraph{Prompt components.} Four semiotic components are factorially
manipulated:
\begin{itemize}
\item \textbf{D} (Domain specification): domain named vs.\ not.
\item \textbf{R} (Role assignment): expert role specified vs.\ not.
\item \textbf{V} (Technical vocabulary): domain terminology included vs.\ absent.
\item \textbf{C} (Explicit constraints): output format, scope, standards stated
      vs.\ absent.
\end{itemize}

Full $2^4 = 16$-condition factorial for three focal domains; four canonical
levels (0--3) for all ten domains. Four canonical levels: \textbf{L0}
$= \{\neg D, \neg R, \neg V, \neg C\}$ (na\"if); \textbf{L1} $= \{D, \neg R,
\neg V, \neg C\}$ (generic); \textbf{L2} $= \{D, R, V, \neg C\}$ (informed);
\textbf{L3} $= \{D, R, V, C\}$ (expert).

\subsubsection{Procedure}

\paragraph{LLM output generation.} $10 \times 16 = 160$ prompts $\times$
$k = 30$ per model $= 14{,}400$ outputs.

\paragraph{Evaluation.} Primary outcome: expert-verified rubric scores.
Two independent domain experts score a stratified subsample of $n \geq 200$
responses (20 per domain per canonical level) against the ground-truth
benchmark on three dimensions: factual accuracy, practical actionability,
domain-specific precision (1--5 scale). Cohen's $\kappa \geq 0.70$ required;
ICC(2,1) reported. LLM-as-judge (GPT-4o) is used as a secondary instrument
\emph{only after} calibration against expert scores ($r \geq 0.75$ required
for acceptance as secondary; otherwise discarded). Note that GPT-4o is
excluded from judging its own outputs to eliminate the most direct
conflict-of-interest; Claude-3.5-Sonnet and Mistral-Large serve as
cross-model judges for GPT-4o outputs.

\paragraph{User study.} $N = 120$ participants recruited via Prolific
Academic, stratified across four education levels $\times$ two domain-familiarity
levels $\times$ three task domains ($n = 10$ per cell). Participants write
one spontaneous prompt per assigned domain (one-shot condition, explicitly
noted as a limitation of ecological validity in Section~\ref{sec:discussion}).
Measured variables:
\begin{itemize}
\item Time on task (seconds; Prolific-recorded)
\item Number of edits before submitting (platform-recorded)
\item Self-reported difficulty (7-point Likert)
\item D/R/V/C component presence coded by two independent raters
      (Cohen's $\kappa \geq 0.80$)
\item Vocabulary sophistication of the prompt (Type-Token Ratio,
      Flesch-Kincaid Grade Level)
\end{itemize}

Participant-authored prompts are fed to GPT-4o; outputs are scored with
the expert rubric. The production-side semiotic cost is operationalized as
the ratio of effort invested (time + edits + self-reported difficulty) to
output quality received: the \emph{effort-to-quality ratio}. The pre-registered
double-bind prediction is that this ratio is \emph{negatively} correlated
with education level and domain familiarity---that is, users with lower
semiotic capital pay more per unit of output quality.

\subsubsection{Analysis}

\paragraph{Output-side.} Mixed-effects model: quality $\sim$ level (ordinal)
$\times$ domain\_complexity + (1$|$domain) + (1$|$participant). Pre-registered
interaction hypothesis: quality gain from L0 to L3 is significantly larger
in high-complexity domains. Monotonicity tested by comparing linear vs.\
threshold (step-function) fits via AIC. Component attribution from the
$2^4$ factorial: V and C are pre-registered as expected primary predictors.

\paragraph{Production-side.} Partial correlation between effort-to-quality
ratio and education level, controlling for domain familiarity. Pre-registered
prediction: negative partial correlation ($r < 0$), indicating that less
educated participants invest more effort per unit of output quality. Moderation
by domain complexity is tested as an exploratory analysis.


%----------------------------------------------------------------------
\section{Implications}
\label{sec:implications}
%----------------------------------------------------------------------

The implications stated below are conditional on experimental confirmation
of the relevant thesis. They are theoretical consequences of the framework,
not reports of results.

\subsection{If Thesis~1 Is Confirmed: Evaluation Should Attend to
Encyclopedic Frames}

The dominant framework for evaluating LLM outputs---the ``hallucination''
framework, measuring deviation from factual reference---is semiotically
inadequate. LLMs do not primarily fail by producing false propositions; they
fail by activating inappropriate encyclopedic pathways. An output that is
``wrong'' is often one that follows a plausible but contextually irrelevant
path through the associative network. Evaluation metrics should therefore
attend not only to factual accuracy but to the appropriateness of the
encyclopedic frame activated. The critical-item accuracy profile from
Experiment~E4, stratified by inferential depth and domain, would provide
a principled basis for such metrics.

\subsection{If Thesis~2 Is Confirmed: Overcoding Is Not Correctable by Debiasing}

If LLM schema conformity exceeds training-corpus baselines (E2), then LLMs
do not merely reproduce the cultural distribution of the training data: they
amplify it. Corrective interventions face a fundamental semiotic obstacle:
they must either restructure the statistical distribution of training data
(raising questions about whose encyclopedia is authorized) or impose post-hoc
constraints on output distributions (substituting one cultural frame for
another, not achieving neutrality). If the diachronic prediction (E5) is
also confirmed---schema conformity increasing across generations for
rising-salience concepts---the implication is that overcoding is a self-reinforcing
dynamic that will intensify as LLM outputs re-enter training corpora. Lotman's
center-periphery dynamic \citep{Lotman1990} is not merely replicated; it is
structurally accelerated.

\subsection{If Thesis~3 Is Confirmed: Semiotic Literacy Is an Educational Priority}

If the double-bind structure is confirmed by the user study (less educated
users invest more effort for worse outputs), the inequality generated by LLMs
is of a specific kind not addressed by access-expansion or domain-knowledge
education alone. It requires cultivating specifically semiotic competences:
the ability to identify what frame is needed, to construct it in language the
model can activate, and to evaluate whether the output has actualized the
intended encyclopedic pathway. This is a new educational challenge that the
semiotic tradition, with its long history of reflection on the conditions of
interpretation, is well positioned to address.

\subsection{Authentication as a Collective Implication of All Three Theses}

The three theses converge on a structural consequence that none entails alone.
If LLM meaning is encyclopedically organized (Thesis~1), then the same
surface discourse activates different pathways depending on contextual
framing---and there is no longer a stable, context-independent sense
\emph{to} be verified. If LLM outputs systematically overcode dominant
cultural frames (Thesis~2), then the cultural authority of a text cannot be
inferred from its fluency or its apparent conformity to convention. If the
quality of LLM output depends on invisible contextual labor (Thesis~3), then
the reader cannot know, from the output alone, whether adequate framing was
supplied. In each case, the provenance of discourse becomes a necessary element
of its interpretation: one must ask not only what a text says but who or what
produced it, by what procedure, under what framing, and within which chain of
accountability.

Authentication is therefore not an external addition to semiosis but, in the
age of LLMs, one of its functional preconditions. This recalls the importance
of paratextual devices in stabilizing meaning \citep{Genette1997}, but the
present situation is more radical: Genette's paratexts orient interpretation;
authentication metadata for LLM-generated content---provenance certificates,
model identifiers, prompt logs, editorial chains---condition whether
interpretation is warranted at all. Developing the taxonomic and normative
vocabulary for these new paratextual functions is a task for semiotics that
no other discipline is currently equipped to perform.


%----------------------------------------------------------------------
\section{Discussion}
\label{sec:discussion}
%----------------------------------------------------------------------

\subsection{What the Experimental Results Can and Cannot Show}

The experiments measure properties of text outputs from statistical models.
They do not directly measure semiotic processes in Eco's theoretical sense,
which are properties of meaning-making in social contexts. The bridging
assumptions in Section~\ref{sec:bridge} specify how computational proxies
are connected to semiotic constructs---but they remain assumptions. Several
implications follow from acknowledging this gap explicitly.

If E4 confirms LLM superiority over the co-occurrence baseline on critical
items, this is strong evidence that context-sensitive disambiguation of the
appropriate kind is occurring. But it does not establish that the mechanism
is encyclopedic organization in Eco's sense of a \emph{culturally weighted}
network; it is consistent with any sufficiently capable contextual language
model. The addition of critical items is designed to maximize the gap between
the co-occurrence and encyclopedic predictions, making the encyclopedic
interpretation the most parsimonious---but not the only possible---explanation
for a positive result.

If E2 confirms that schema conformity for charged concepts exceeds the
corpus baseline, the model-amplification interpretation is the most direct
one: the LLM has learned not merely to reproduce the training distribution
but to concentrate it onto the dominant cultural prototype. This is a
meaningful finding independently of whether it is called ``overcoding''---
though the semiotic vocabulary adds precision by specifying the social mechanism
(cultural code superimposition) that the statistical phenomenon instantiates.

What the experiments cannot show is whether these properties generalize beyond
instruction-tuned English-language LLMs, beyond the concepts and domains
investigated, or beyond the synchronic moment at which data are collected.
Scope is explicitly restricted to instruction-tuned English LLMs; claims for
base models, non-English systems, or future model architectures are explicitly
withheld pending further investigation.

\subsection{The Intentionality Question}

The claim that LLMs generate discourse ``without intentional consciousness''
(abstract) positions this work within a contested philosophical terrain.
The Chinese Room argument \citep{Searle1980} holds that symbol manipulation
without grounding cannot constitute genuine understanding; the symbol grounding
problem \citep{Harnad1990} argues that ungrounded symbols lack original meaning.
These positions imply that LLM semiosis is derived or simulated rather than
genuine. The present article takes no position on the grounding question at
the level of individual signs. What it does claim is that, at the level of
\emph{social function}, LLM outputs circulate as semiotic objects: they are
interpreted, acted upon, cited, and credited in ways that produce real
social effects regardless of their internal grounding status. Semiosis is
here a property of a social process, not of a cognitive state. Eco's own
framework, which treats semiosis as a property of sign-systems rather than
of sign-users, supports this social-functional approach
\citep{Eco1976,Eco1979}.

\subsection{Model Dependency and Scope}

Claims are restricted to patterns replicating across all three models. Where
results diverge, the most likely explanatory variable is the RLHF and
safety-tuning procedure, which differs substantially across providers. A
finding that GPT-4o shows lower entropy than Llama-3 on culturally charged
concepts, for example, may reflect differential safety filtering rather than
a difference in structural overcoding tendency. Such divergences are noted and
analyzed but not treated as counterevidence against the class-level claim.

\subsection{Ecological Validity of the User Study}

The user study in E3 employs a one-shot prompting condition, which does not
capture the iterative, goal-directed character of naturalistic LLM use. Real
users revise prompts across multiple turns; their semiotic cost is distributed
across a session, not concentrated in a single composition event. The one-shot
design is a necessary simplification for experimental control, but it likely
\emph{underestimates} the total semiotic cost, particularly for users with
lower competence who iterate more. Claims about the magnitude of the effort-
to-quality differential are therefore conservative; the directional claim
is unaffected.

\subsection{Implications for Computational Fairness and Policy}

The overcoding thesis implies that ``debiasing'' LLMs is semiotically more
complex than its typical presentation suggests. If schema conformity in LLM
outputs exceeds training-corpus levels (E2 confirms model amplification), then
debiasing cannot be achieved by curating the training data alone: the amplification
is a property of the generative mechanism, not only of the input distribution.
Interventions at the output level risk substituting one dominant frame for
another without addressing the structural tendency toward overcoding. Effective
intervention requires either architectural changes that reduce conformity-amplification,
or explicit diversity-injection at the generation level, with full awareness
that any such injection is itself a culturally motivated choice.

The semiotic cost thesis implies that the deployment of LLMs in high-stakes
domains---legal advice, medical information, administrative communication---may
exacerbate existing inequalities unless accompanied by semiotic literacy
education. Access to the tool is a necessary but insufficient condition for
equitable benefit; what is additionally required is the specific competence
to frame requests in ways that activate the model's relevant encyclopedic
pathways. Developing this competence requires educational programmes that
explicitly address the pragmatics of machine-directed communication---a new
form of literacy that combines elements of rhetoric, domain knowledge, and
semiotic theory.


%----------------------------------------------------------------------
\section{Conclusion}
\label{sec:conclusion}
%----------------------------------------------------------------------

LLMs do not merely add a new instrument to the repertoire of communication.
They modify the conditions under which discourse is produced, contextualized,
validated, and trusted. They industrialize the generation of plausible
signifying sequences in the specific sense that separates textual form from
the labor of meaning: without the intentional positioning of an author,
without accountability to an interpretive community, and without the
friction of communicative failure that ordinarily regulates the social
circulation of discourse.

Eco's semiotics is particularly well equipped to analyze this transformation
because it developed precisely the concepts that describe what LLMs exploit:
the encyclopedia (the cultural network that makes plausibility possible),
overcoding (the mechanism by which semiosis favors dominant convention), and
interpretive cooperation (the pragmatic labor that grounding-free discourse
transfers onto the reader). The three theses advanced in this article each
operationalize one of these concepts as a falsifiable empirical claim, and the
five-experiment design tests each claim at two levels---mechanistic and
distributional for Thesis~1, synchronic and diachronic for Thesis~2, output-side
and production-side for Thesis~3.

The explicit theoretical-empirical bridge (Section~\ref{sec:bridge}) specifies
what each proxy measures, what the bridging assumption is, and under what
conditions the semiotic interpretation would be disconfirmed. Pre-registered
adjudication rules specify how conflicting evidence within a thesis is to be
interpreted. This level of specification is unusual in semiotic research, but
it is the minimum required for claims about LLM behavior---systems whose
outputs are directly observable and whose statistical structure is measurable---
to carry empirical weight.

If the experimental results confirm these theses, then semiotics has a
distinctive contribution to make to the evaluation, governance, and education
surrounding LLMs: not as an ornamental commentary but as one of the disciplines
best equipped to explain what is most consequential about them. If the results
disconfirm them, the theoretical framework must be revised---and the revision
will itself constitute a contribution to semiotic theory.


%----------------------------------------------------------------------
% References
%----------------------------------------------------------------------

\bigskip
\begin{thebibliography}{99}

\bibitem[Adorno and Horkheimer(1944)]{AdornoHorkheimer1944}
Adorno, T.\,W., \& Horkheimer, M. (1944).
\newblock \emph{Dialectic of Enlightenment}.
\newblock New York: Herder and Herder (1972 English translation).

\bibitem[Barthes(1977)]{Barthes1977}
Barthes, R. (1977).
\newblock The Death of the Author.
\newblock In \emph{Image-Music-Text} (pp.~142--148). London: Fontana.

\bibitem[Bateman(2024)]{Bateman2024}
Bateman, J.\,A. (2024).
\newblock What Are Large Language Models Doing? A Semiotic Perspective.
\newblock Research Seminar in Language Technology, University of Helsinki,
16 May 2024.
\newblock \url{https://blogs.helsinki.fi/language-technology/files/2024/05/bateman-semiotic-llm-2024.pdf}

\bibitem[Bender et al.(2021)]{Bender2021}
Bender, E.\,M., Gebru, T., McMillan-Major, A., \& Shmitchell, S. (2021).
\newblock On the Dangers of Stochastic Parrots: Can Language Models Be Too Big?
\newblock In \emph{Proceedings of FAccT~2021} (pp.~610--623). ACM.
\newblock \url{https://doi.org/10.1145/3442188.3445922}

\bibitem[Bevilacqua and Navigli(2020)]{Bevilacqua2020}
Bevilacqua, M., \& Navigli, R. (2020).
\newblock Breaking Through the 80\% Glass Ceiling: Raising the State of the
Art in Word Sense Disambiguation by Incorporating Knowledge Graph Information.
\newblock In \emph{Proceedings of ACL~2020} (pp.~2854--2864).
\newblock \url{https://doi.org/10.18653/v1/2020.acl-main.255}

\bibitem[Bourdieu(1991)]{Bourdieu1991}
Bourdieu, P. (1991).
\newblock \emph{Language and Symbolic Power}.
\newblock Cambridge: Polity Press.

\bibitem[Eco(1976)]{Eco1976}
Eco, U. (1976).
\newblock \emph{A Theory of Semiotics}.
\newblock Bloomington: Indiana University Press.

\bibitem[Eco(1979)]{Eco1979}
Eco, U. (1979).
\newblock \emph{The Role of the Reader: Explorations in the Semiotics of Texts}.
\newblock Bloomington: Indiana University Press.

\bibitem[Eco(1984)]{Eco1984}
Eco, U. (1984).
\newblock \emph{Semiotics and the Philosophy of Language}.
\newblock Bloomington: Indiana University Press.

\bibitem[Eco(1990)]{Eco1990}
Eco, U. (1990).
\newblock \emph{The Limits of Interpretation}.
\newblock Bloomington: Indiana University Press.

\bibitem[Fiske et al.(2002)]{Fiske2002}
Fiske, S.\,T., Cuddy, A.\,J.\,C., Glick, P., \& Xu, J. (2002).
\newblock A Model of (Often Mixed) Stereotype Content: Competence and Warmth
Respectively Follow from Perceived Status and Competition.
\newblock \emph{Journal of Personality and Social Psychology}, 82(6), 878--902.
\newblock \url{https://doi.org/10.1037/0022-3514.82.6.878}

\bibitem[Genette(1997)]{Genette1997}
Genette, G. (1997).
\newblock \emph{Paratexts: Thresholds of Interpretation}.
\newblock Cambridge: Cambridge University Press.
\newblock \url{https://doi.org/10.1017/CBO9780511549373}

\bibitem[Hall of Mirrors(2026)]{HallMirrors2026}
(2026).
\newblock Language Models' Hall of Mirrors Problem: Why AI Alignment Requires
Peircean Semiosis.
\newblock \emph{Philosophy \& Technology}, 38.
\newblock \url{https://link.springer.com/article/10.1007/s13347-025-00975-5}

\bibitem[Harnad(1990)]{Harnad1990}
Harnad, S. (1990).
\newblock The Symbol Grounding Problem.
\newblock \emph{Physica D: Nonlinear Phenomena}, 42(1--3), 335--346.
\newblock \url{https://doi.org/10.1016/0167-2789(90)90087-6}

\bibitem[Lackov\'a Bennett(2025)]{LackovaBennett2025}
Lackov\'a Bennett, \v{L}. (2025).
\newblock Toward a Biosemiotics Framework for AI: Folding and the Dynamics
of Meaning.
\newblock \emph{Sign Systems Studies}, 53(3--4), 444--463.
\newblock \url{https://doi.org/10.12697/SSS.2025.53.3-4.07}

\bibitem[Leone(2023)]{Leone2023}
Leone, M. (2023).
\newblock The Main Tasks of a Semiotics of Artificial Intelligence.
\newblock \emph{Language and Semiotic Studies}, 9(1), 1--13.
\newblock \url{https://doi.org/10.1515/lass-2022-0006}

\bibitem[Lotman(1990)]{Lotman1990}
Lotman, Y.\,M. (1990).
\newblock \emph{Universe of the Mind: A Semiotic Theory of Culture}.
\newblock Bloomington: Indiana University Press.

\bibitem[LSM(2025)]{LSM2025}
(2025).
\newblock Beyond Tokens: Introducing Large Semiosis Models (LSMs) for Grounded
Meaning in Artificial Intelligence.
\newblock Preprints.org, 202504.1830.
\newblock \url{https://www.preprints.org/manuscript/202504.1830}

\bibitem[Miller(1995)]{Miller1995}
Miller, G.\,A. (1995).
\newblock WordNet: A Lexical Database for English.
\newblock \emph{Communications of the ACM}, 38(11), 39--41.
\newblock \url{https://doi.org/10.1145/219717.219748}

\bibitem[Not Minds(2025)]{NotMinds2025}
(2025).
\newblock Not Minds, but Signs: Reframing LLMs through Semiotics.
\newblock arXiv:2505.17080.
\newblock \url{https://arxiv.org/abs/2505.17080}

\bibitem[Nunes and Antunes(2024)]{NunesAntunes2024}
Nunes, D., \& Antunes, L. (2024).
\newblock Machines of Meaning.
\newblock arXiv:2412.07975.
\newblock \url{https://arxiv.org/abs/2412.07975}

\bibitem[Peirce(1931)]{Peirce1931}
Peirce, C.\,S. (1931--1958).
\newblock \emph{Collected Papers of Charles Sanders Peirce}, Vols.~1--8.
\newblock Cambridge, MA: Harvard University Press.

\bibitem[Pennington et al.(2014)]{Pennington2014}
Pennington, J., Socher, R., \& Manning, C.\,D. (2014).
\newblock GloVe: Global Vectors for Word Representation.
\newblock In \emph{Proceedings of EMNLP~2014} (pp.~1532--1543).
\newblock \url{https://doi.org/10.3115/v1/D14-1162}

\bibitem[Piantadosi and Hill(2022)]{PiantadosiHill2022}
Piantadosi, S.\,T., \& Hill, F. (2022).
\newblock Meaning without Reference in Large Language Models.
\newblock arXiv:2208.02957.
\newblock \url{https://arxiv.org/abs/2208.02957}

\bibitem[Popper(1963)]{Popper1963}
Popper, K.\,R. (1963).
\newblock \emph{Conjectures and Refutations: The Growth of Scientific Knowledge}.
\newblock London: Routledge.

\bibitem[Searle(1980)]{Searle1980}
Searle, J.\,R. (1980).
\newblock Minds, Brains, and Programs.
\newblock \emph{Behavioral and Brain Sciences}, 3(3), 417--424.
\newblock \url{https://doi.org/10.1017/S0140525X00005756}

\bibitem[Signata(2025)]{Signata2025}
\emph{Signata: Annales des S\'emiotiques} (2025).
\newblock Special Issue: Semiotics of Space in the Era of Computational Logic
and Digital Practices.
\newblock \url{https://journals.openedition.org/signata/5880}

\bibitem[Tomalin(2025)]{Tomalin2025}
Tomalin, M. (2025).
\newblock Multimodal Social Semiotics and the Challenge of Artificial Intelligence.
\newblock \emph{Visual Communication}.
\newblock \url{https://doi.org/10.1177/26349795251327939}

\bibitem[Vaswani et al.(2017)]{Vaswani2017}
Vaswani, A., Shazeer, N., Parmar, N., Uszkoreit, J., Jones, L., Gomez, A.\,N.,
Kaiser, \L., \& Polosukhin, I. (2017).
\newblock Attention Is All You Need.
\newblock In \emph{Advances in Neural Information Processing Systems}, 30.
\newblock \url{https://doi.org/10.48550/arXiv.1706.03762}

\bibitem[Vromen(2024)]{Vromen2024}
Vromen, E. (2024).
\newblock Language Models as Semiotic Machines: Reconceptualizing AI Language
Systems through Structuralist and Post-Structuralist Theories of Language.
\newblock arXiv:2410.13065.
\newblock \url{https://arxiv.org/abs/2410.13065}

\bibitem[Wang et al.(2022)]{Wang2022}
Wang, L., Yang, N., Huang, X., Jiao, B., Yang, L., Jiang, D., Majumder, R.,
\& Wei, F. (2022).
\newblock Text Embeddings by Weakly-Supervised Contrastive Pre-Training.
\newblock arXiv:2212.03533.
\newblock \url{https://arxiv.org/abs/2212.03533}

\bibitem[Warschauer(2004)]{Warschauer2004}
Warschauer, M. (2004).
\newblock \emph{Technology and Social Inclusion: Rethinking the Digital Divide}.
\newblock Cambridge, MA: MIT Press.

\end{thebibliography}

\end{document}
